\item \model{GLM-5}

\paragraph*{مقدمه و پارادایم داده‌های پس‌آموزش}
فاز پس‌آموزش (\en{Post-Training}) در مدل \model{GLM-5} با هدف گذار ساختاری از پارادایم «کدنویسی حسی» (\en{Vibe Coding})—که مبتنی بر پرامپت‌نویسی انسانی است—به سمت «مهندسی عامل‌محور» (\en{Agentic Engineering}) بازطراحی شده است؛ جایی که مدل به عنوان یک عامل خودکار، وظایف برنامه‌ریزی، پیاده‌سازی، اجرا و خوداصلاحی را در افق‌های زمانی طولانی بر عهده می‌گیرد. برای این منظور، خط لوله داده‌های پس‌آموزش با بسط طول بافت تا سقف \textbf{۲۰۲,۷۵۲ توکن} در قالب یک زنجیره چندمرحله‌ای شامل تنظیم دقیق با‌نظارت (\en{SFT})، یادگیری تقویتی استدلال (\en{Reasoning RL})، یادگیری تقویتی عامل‌محور (\en{Agentic RL})، یادگیری تقویتی عمومی (\en{General RL}) و تقطیر هم‌خط بین‌مرحله‌ای (\en{On-Policy Cross-Stage Distillation}) سازمان‌دهی شده است.

\paragraph*{ترکیب‌بندی و پالایش پیکره داده‌های تنظیم دقیق بانظارت (\en{SFT Data})}
پیکره داده‌های \en{SFT} در \model{GLM-5} نسبت به نسل‌های پیشین به شکل چشمگیری در حوزه‌های کدنویسی و تعاملات عاملی گسترش یافته و در سه محور ساختار یافته است:
\begin{itemize}
    \item \textbf{گفت‌وگوی عمومی (\en{General Chat}):} شامل داده‌های پرسش و پاسخ، نگارش، ترجمه، تعاملات طولانی و بازی نقش (\en{Role-playing}) چندزبانه. داده‌های بازی نقش بر پایه پنج بعد کیفی (پیروی از دستور، بیان‌گری زبانی، خلاقیت، انسجام منطقی و پایداری در مکالمات طولانی) توسط ارزیاب‌های خودکار و متخصصان انسانی پالایش شدند. سبک پاسخ‌ها نیز به سوی ساختاری منطقی‌تر و موجزتر بهینه‌سازی شد.
    \item \textbf{داده‌های استدلال تحلیلی (\en{Reasoning Data}):} در حوزه استدلال منطقی، داده‌ها از مسائل قابل راستی‌آزمایی عینی با روش نمونه‌برداری حذفی (\en{Rejection Sampling}) استخراج شدند. در مسائل ریاضی و علوم، فیلتر مبتنی بر دشواری (\en{Difficulty-based Filtering}) اعمال گردید تا تنها مسائلی حفظ شوند که مدل \model{GLM-4.7} توانایی حل آن‌ها را نداشته است.
    \item \textbf{داده‌های کدنویسی و عاملی (\en{Coding \& Agent Data}):} شامل کدهای فرانت‌اند و بک‌اند، فراخوانی ابزارها و ردپای عامل‌های حل مسأله در محیط‌های واقعی. داده‌ها با الگوریتم‌های یادگیری تقویتی خبره و نمونه‌برداری حذفی استخراج شدند. بخش‌های اشتباه در مسیر اجرای عامل حذف نشدند، بلکه در تابع زیان آموزش ماسک‌گذاری شدند (\en{Masked out}):
    \begin{equation}
        \mathcal{L}_{\text{SFT}}(\theta) = - \sum_{t \in \mathcal{T}_{\text{valid}}} \log \pi_\theta(y_t \mid x, y_{<t})
    \end{equation}
    این رویکرد به مدل می‌آموزد که الگوهای تصحیح خطا (\en{Error Correction}) را بدون بازتولید یا تقویت اقدامات غلط فراگیرد.
\end{itemize}

\paragraph*{قالب‌بندی داده‌ها و حالات سه‌گانه تفکر (\en{Thinking Modes})}
قالب پرامپت‌ها در داده‌های پس‌آموزش به‌گونه‌ای مهندسی شده است که مدل را به سه الگوی رفتاری ساختاریافته مجهز می‌سازد:
\begin{itemize}
    \item \textbf{تفکر درهم‌تنیده (\en{Interleaved Thinking}):} تعبیه بلوک‌های استدلال پیش از هر پاسخ متنی و پیش از هر فراخوانی ابزار (\en{Tool Call}) جهت بهبود دقت عملیاتی.
    \item \textbf{تفکر حفظ‌شده (\en{Preserved Thinking}):} در تعاملات چندنوبته مهندسی نرم‌افزار، زنجیره تفکر نوبت‌های پیشین در زمینه ذخیره می‌شود تا از استدلال مجدد از نقطه صفر و ایجاد تناقض‌های منطقی در وظایف طولانی جلوگیری شود.
    \item \textbf{تفکر در سطح نوبت (\en{Turn-level Thinking}):} قابلیت تنظیم فعال/غیرفعال بودن تفکر در هر نوبت بر اساس میزان پیچیدگی درخواست به منظور کاهش تأخیر زمانی.
\end{itemize}

\paragraph*{داده‌ها و فرمولاسیون یادگیری تقویتی استدلال (\en{Reasoning RL})}
در مرحله استدلال تحلیلی، داده‌ها در چهار حوزه متعادل توزیع شده‌اند:
\begin{itemize}
    \item \textbf{ریاضیات و علوم:} گردآوری از دادگان متن‌باز معتبر نظیر \en{Nemotron-Math} \cite{nemotronmath2025} و راهکار برنده رقابت \en{AIMO-2} با پایگاه داده \en{OpenMathReasoning} \cite{moshkov2025aimo2}. مسائل به گونه‌ای فیلتر شدند که برای \model{GLM-4.7} غیرقابل حل بوده، اما معلمان قدرتمند نظیر \en{GPT-5.2 (xhigh)} \cite{openai2025gpt52} و \en{Gemini 3 Pro Preview} \cite{deepmind2025gemini3} توانایی حل گام‌به‌گام آن‌ها را داشته باشند.
    \item \textbf{کدنویسی رقابتی و علمی:} مسائل مسابقات برنامه‌نویسی از پلتفرم \en{Codeforces}، مجموعه‌داده‌های \en{TACO} \cite{li2023taco} و \en{SYNTHETIC-2-RL} \cite{primeintellect2025synthetic2}.
    \item \textbf{استدلال یکپارچه با ابزار (\en{Tool-Integrated Reasoning - TIR}):} نمونه مسائل ریاضی و مهندسی که مستلزم فراخوانی مفسر پایتون و ابزارهای محاسباتی بیرونی هستند.
\end{itemize}

آموزش بر پایه نسخه اصلاح‌شده الگوریتم \en{GRPO} \cite{shao2024deepseekmath} و تکنیک \en{IcePop} \cite{zhao2025icepop} برای حذف واگرایی میان سیاست استنتاج و سیاست گرادیان صورت گرفته است:
\begin{equation}
\begin{aligned}
    \mathcal{L}(\theta) = -\mathbb{E}_{x \sim \mathcal{D}, \{y_i\}_{i=1}^G \sim \pi_{\theta_{\text{old}}}^{\text{infer}}} \Bigg[ & \frac{1}{G} \sum_{i=1}^G \frac{1}{|y_i|} \sum_{t=1}^{|y_i|} \text{pop}(\rho_{i,t}, 1/\beta, \beta) \\
    & \cdot \min \left( r_{i,t}\hat{A}_{i,t}, \, \text{clip}(r_{i,t}, 1-\epsilon_{\text{low}}, 1+\epsilon_{\text{high}})\hat{A}_{i,t} \right) \Bigg]
\end{aligned}
\end{equation}
که در آن نسبت عدم‌انطباق استنتاج-آموزش و عملگر سرکوب خطای $\text{pop}$ عبارتند از:
\begin{equation}
    \rho_{i,t} = \frac{\pi_{\theta_{\text{old}}}^{\text{train}}(y_{i,t} \mid x, y_{i,<t})}{\pi_{\theta_{\text{old}}}^{\text{infer}}(y_{i,t} \mid x, y_{i,<t})}, \quad
    \text{pop}(\rho, 1/\beta, \beta) = \begin{cases} \rho, & \text{if } 1/\beta \le \rho \le \beta \\ 0, & \text{otherwise} \end{cases}
\end{equation}
آموزش این بخش با اندازه گروه $G=32$، اندازه دسته ۳۲ و هایپرپارامترهای $\beta=2$، $\epsilon_{\text{low}}=0.2$ و $\epsilon_{\text{high}}=0.28$ انجام شده است.

\paragraph*{محیط‌های اجرایی و ساخت دادگان عامل‌محور (\en{Agentic RL Environments})}
برای آموزش مدل در نقش عامل مهندسی، داده‌ها از متن‌های ایستا فراتر رفته و در بستر بیش از ۱۰ هزار محیط اجرایی ابطال‌پذیر تعاملی تولید شدند:
\begin{itemize}
    \item \textbf{محیط‌های مهندسی نرم‌افزار (\en{SWE Environments}):} با اقتباس از چارچوب \en{RepoLaunch} \cite{zhang2025swebenchlive}، بیش از \textbf{۱۰,۰۰۰ محیط اجرایی} بر پایه مخازن گیت‌هاب در ۹ زبان برنامه‌نویسی (\en{Python}، \en{Java}، \en{Go}، \en{C}، \en{C++}، \en{JavaScript}، \en{TypeScript}، \en{PHP} و \en{Ruby}) ایجاد شد. آزمون‌های شکست-به-موفقیت (\en{Fail-to-Pass - F2P}) و پایداری (\en{Pass-to-Pass - P2P}) توسط لاگ‌خوان‌های مبتنی بر مدل زبانی استخراج شدند.
    \item \textbf{محیط‌های ترمینال (\en{Terminal Environments}):} تولید داده‌ها در ساختار استاندارد کانتینری \en{Harbor} \cite{harbor2026} انجام گرفت. این فرآیند از دو مسیر انجام شد: ۱) سنتز سه‌مرحله‌ای از وظایف بذر مهندسی نرم‌افزار با دقت ساخت داکر بیش از ۹۰٪، ۲) استخراج خودکار و حلقه-بسته از صفحات وب فنی که در آن عامل با اجرای اسکریپت‌های اعتبارسنجی \en{Harbor} به صورت خودکار وظیفه تولیدی را اصلاح و نهایی می‌کرد.
    \item \textbf{وظایف جستجوی عمیق چندمرحله‌ای (\en{Deep-Search Tasks}):} پایگاه داده‌ای از گراف دانش وب (\en{WKG}) بر پایه ۲ میلیون صفحه با چگالی اطلاعاتی بالا ساخته شد. سوالات از روی زنجیره‌های چندموجودیتی استخراج و از یک فیلتر سه‌مرحله‌ای عبور داده شدند: ۱) حذف سوالات قابل حل بدون ابزار (در ۱ تلاش از ۸ اجرا)، ۲) حذف سوالات قابل حل با جستجوی تک‌مرحله‌ای، ۳) اعتبارسنجی دوطرفه تطابق شواهد وب با پاسخ مرجع توسط عامل داور.
    \item \textbf{محیط تولید اسلاید متنی-بصری (\en{Slide Generation}):} مدلسازی داده‌ها بر پایه ساختار \en{HTML} با پاداش سه‌سطحی: سطح ۱ (ویژگی‌های نشانه‌گذاری استایل، رنگ و حذف تصاویر توهمی)، سطح ۲ (ویژگی‌های رندرینگ زمان اجرا در \en{DOM} نظیر ابعاد و جلوگیری از هک پاداش با قطع نامتقارن متن)، و سطح ۳ (ارزیابی بصری توزیع فضاهای خالی). در این فرایند انطباق ساختاری با نسبت ابعاد ۱۶:۹ از ۴۰٪ به \textbf{۹۲٪} افزایش یافت.
\end{itemize}

\paragraph*{پایدارسازی آماری داده‌ها در یادگیری تقویتی ناهمگام (\en{Asynchronous Stability})}
تولید برون‌خط و هم‌زمان مسیرهای اجرایی عامل‌ها نیازمند کنترل آماری عدم‌انطباق سیاست‌ها بود:
\begin{itemize}
    \item \textbf{دروازه توکن‌به‌توکن (\en{TITO Gateway}):} ثبت مستقیم شناسه‌های توکن و فراداده‌ها در خروجی موتور استنتاج به منظور رفع خطاهای ناشی از توکنایزیشن مجدد در سمت آموزش.
    \item \textbf{نمونه‌برداری اهمیتی دوطرفه مستقیم (\en{Direct Double-sided Importance Sampling}):} کنترل بازه گرادیان در محدوده $[1-\epsilon_\ell, 1+\epsilon_h]$ بدون نیاز به فراخوانی پرهزینه سیاست قدیمی:
    \begin{equation}
        \mathcal{L}(\theta) = \mathbb{E}_t \left[ f(r_t(\theta), \epsilon_\ell, \epsilon_h) \hat{A}_t \log \pi_\theta(a_t \mid s_t) \right]
    \end{equation}
    \begin{equation}
        r_t(\theta) = \exp \left( \log \pi_\theta(a_t \mid s_t) - \log \pi_{\text{rollout}}(a_t \mid s_t) \right)
    \end{equation}
    \begin{equation}
        f(x; \epsilon_\ell, \epsilon_h) = \begin{cases} x, & \text{if } 1 - \epsilon_\ell < x < 1 + \epsilon_h \\ 0, & \text{otherwise} \end{cases}
    \end{equation}
    \item \textbf{پالایش داده‌های بیات و خطاهای محیطی:} حذف مسیرهایی که اختلاف نسخه مدل تولیدکننده آن‌ها با نسخه جاری آموزش بیش از حد آستانه بود ($w' - w_0 > \tau$). همچنین حذف داده‌های حاصل از فروپاشی کانتینر داکر و تکمیل داده‌های گروهی (\en{Padding}) در صورت وجود حداقل نیمی از نمونه‌های معتبر.
\end{itemize}

\paragraph*{یادگیری تقویتی عمومی و تقطیر هم‌خط بین‌مرحله‌ای}
در فاز یادگیری عمومی، اهداف بهینه‌سازی در سه بعد صحت بنیادین (کاهش توهم و خطاهای ساختاری)، هوش هیجانی (همدلی و گفت‌وگوی روان) و کیفیت تخصصی وظایف تنظیم شدند. برای نظارت بر این ابعاد، سیستمی ترکیبی از پاداش‌های قانون‌محور، مدل‌های پاداش پیامد (\en{ORMs}) و مدل‌های پاداش زایندگی (\en{GRMs}) به کار رفت. افزون بر این، به منظور ممانعت از تولید خروجی‌های بیش‌ازحد کلیشه‌ای و ماشینی، پاسخ‌های تدوین‌شده توسط متخصصان برجسته انسانی (\en{Human-authored Responses}) به عنوان لنگرهای کیفی سبک وارد جریان آموزش شدند.

در نهایت، برای رفع افت تدریجی توانمندی‌های پیشین (\en{Catastrophic Forgetting})، تقطیر هم‌خط بین‌مرحله‌ای بر روی توزیع داده‌های مراحل استدلال و عمومی با اندازه دسته ۱۰۲۴ و اندازه گروه $G=1$ اجرا شد:
\begin{equation}
    \hat{A}_{i,t} = \text{sg} \left[ \log \frac{\pi_{\theta_{\text{teacher}}}^{\text{infer}}(y_{i,t} \mid x, y_{i,<t})}{\pi_\theta^{\text{train}}(y_{i,t} \mid x, y_{i,<t})} \right]
\end{equation}

\paragraph*{ارزیابی‌های تجربی روی مجموعه‌داده‌های عامل‌محور و دنیای واقعی}
به منظور ارزیابی کارایی داده‌های پس‌آموزش در شرایط واقعی مهندسی، بنچمارک داخلی جامع \en{CC-Bench-V2} طراحی گردید. آمار ساختار داده‌ها و نتایج عملکردی مدل در جدول‌های زیر خلاصه شده است.

\begin{table}[htbp]
\centering
\caption{توزیع سناریوهای کاربردی فرانت‌اند و واحدهای اعتبارسنجی در \en{CC-Bench-V2}}
\label{tab:glm5_frontend_dist}
\resizebox{\textwidth}{!}{%
\begin{tabular}{llcc}
\toprule
\textbf{دسته‌بندی} & \textbf{شرح سناریو} & \textbf{تعداد وظایف (\en{\# Tasks})} & \textbf{تعداد چک‌لیست‌ها (\en{\# Checkitems})} \\
\midrule
سامانه‌های سازمانی (\en{Business Systems}) & مدیریت فرآیندها و داده‌های تجاری & 42 & 167 \\
بازی‌های تعاملی وب (\en{Web Games}) & تعاملات گرافیکی و برنامه‌های سرگرمی & 40 & 163 \\
رندرینگ برداری (\en{SVG/Canvas}) & ترسیمات پویا و مصورسازی‌های تعاملی & 32 & 166 \\
ابزارهای خلاقانه (\en{Creative Tools}) & ابزارهای ویرایش و تولید محتوای برخط & 28 & 160 \\
صفحات معرفی (\en{Showcase Pages}) & بازنمایی اطلاعات و طراحی صفحات فرود & 27 & 115 \\
فرم‌ها و جداول (\en{Forms \& Tables}) & ساختارهای ورود داده و اعتبارسنجی ورودی & 26 & 93 \\
مصورسازی داده (\en{Data Visualization}) & نمودارهای تحلیلی و گرافیکی وب & 25 & 85 \\
\midrule
\textbf{مجموع کل وظایف فرانت‌اند} & پوشش کامل پشته‌های \en{HTML/JS} (۱۱۳)، \en{React} (۵۸) و \en{Vue} (۴۹) & \textbf{220} & \textbf{969} \\
\bottomrule
\end{tabular}%
}
\end{table}

\begin{table}[htbp]
\centering
\caption{ارزیابی عملکرد مدل در بخش‌های فرانت‌اند، بک‌اند و افق طولانی بنچمارک \en{CC-Bench-V2}}
\label{tab:glm5_ccbench_eval}
\resizebox{\textwidth}{!}{%
\begin{tabular}{lllccc}
\toprule
\textbf{حوزه مهندسی} & \textbf{پشته یا نوع وظیفه} & \textbf{متریک ارزیابی} & \textbf{\model{GLM-5}} & \textbf{\model{GLM-4.7}} & \textbf{\en{Claude Opus 4.5}} \\
\midrule
\multirow{3}{*}{\textbf{فرانت‌اند} (\en{Frontend})} 
& \en{HTML/CSS/JS} & \en{ISR / CSR} (\%) & \textbf{38.9 / 76.3} & 35.4 / 64.9 & 52.2 / 82.2 \\
& \en{React} & \en{ISR / CSR} (\%) & \textbf{34.6 / 71.0} & 17.2 / 49.4 & 39.7 / 70.7 \\
& \en{Vue} & \en{ISR / CSR} (\%) & \textbf{32.7 / 77.1} & 24.5 / 53.8 & 46.9 / 74.3 \\
\midrule
\textbf{صحت ساخت} (\en{Build}) & \en{React / Vue / Svelte / Next} & \en{BSR} (\%) & \textbf{100 / 100 / 100 / 95.0} & 65.0 / 70.0 / 60.0 / 70.0 & 95.0 / 100 / 90.0 / 80.0 \\
\midrule
\textbf{بک‌اند} (\en{Backend}) & مهندسی نرم‌افزار (۸۵ تسک در ۶ زبان) & \en{Pass@1} (\%) & \textbf{25.8} & 19.6 & 26.9 \\
\midrule
\multirow{2}{*}{\textbf{افق طولانی} (\en{Long-horizon})} 
& ناوبری مخازن بزرگ (\en{Repo Exploration}) & \en{Pass@1} (\%) & \textbf{65.6} & 47.8 & 64.5 \\
& وظایف زنجیره‌ای متوالی (\en{Chained Tasks}) & \en{Pass@1} (\%) & \textbf{52.3} & 43.0 & 61.6 \\
\bottomrule
\end{tabular}%
}
\end{table}

همچنین مدل \model{GLM-5} بر روی مجموعه‌داده‌های عملیاتی ارزیابی گردید که دستاوردهای زیر حاصل شد:
\begin{itemize}
    \item \textbf{ترجمه ماشینی واقعی (\en{ZMultiTransBench}):} شامل ۱,۲۲۰ نمونه از سناریوهای متداول کاربری در ۷ جفت‌زبان از مبدا چینی به اسپانیایی (۳۰۰)، روسی (۲۵۰)، فرانسوی (۲۲۰)، کره‌ای (۲۰۰)، ژاپنی (۱۵۰)، عربی (۵۰) و آلمانی (۵۰). امتیاز الو مدل از ۱۰۱۶ در \model{GLM-4.7} به \textbf{۱۰۵۰} در \model{GLM-5} افزایش یافت.
    \item \textbf{ترجمه زبانی دشوار (\en{MENT-SNS} \cite{tian2026beyond}):} ۷۵۳ جفت‌جمله انگلیسی-چینی دارای کنایه و استعاره؛ امتیاز الو از ۹۹۳ به \textbf{۱۰۱۳} ارتقا یافت.
    \item \textbf{گفت‌وگوی چندزبانه (\en{ZMultiDialBench}):} ۱۴۱ سناریوی چندزبانه ارزیابی‌شده توسط متخصصان؛ نمره میانگین از ۸.۴۴ به \textbf{۸.۵۷} بهبود یافت.
    \item \textbf{پیروی از دستورات پیچیده (\en{IF-Badcase}):} ۴۵۰ نمونه پرامپت با محدودیت‌های چندگانه استخراج‌شده از محیط استقرار؛ نرخ موفقیت از ۷۸.۵٪ به \textbf{۸۳.۲٪} رسید.
    \item \textbf{فراخوانی ابزار (\en{ToolCall-Badcase}):} ۲۰۰ نمونه استخراج‌شده از موارد شکست کاربران در محیط عملیاتی؛ دقت تطابق ساختار و آرگومان‌ها با جهشی خیره‌کننده از ۶۰.۸٪ به \textbf{۹۵.۸٪} رسید.
\end{itemize}